% Part: first-order-logic
% Chapter: syntax-and-semantics
% Section: covered-structures

\documentclass[../../../include/open-logic-section]{subfiles}

\begin{document}

\olfileid{fol}{syn}{cov}

\olsection{一阶语言的覆盖结构}

\begin{explain}
回想一下，如果一个项不包含任何变元，则称其为\emph{闭项}。
\end{explain}

\begin{defn}[闭项的值]
若 $t$ 是语言~$\Lang L$ 的闭项，且 $\Struct M$ 是~$\Lang L$ 的结构，则 $t$ 的\emph{值}~$\Value{t}{M}$ 定义如下：
\begin{enumerate}
\item 若 $t$ 仅为常项符号~$c$，则 $\Value{c}{M} = \Assign{c}{M}$。
\item 若 $t$ 形如 $\Atom{f}{t_1, \ldots, t_n}$，则
  \[
  \Value{t}{M} = \Assign{f}{M}(\Value{t_1}{M}, \ldots,
  \Value{t_n}{M}).
  \]
\end{enumerate}
\end{defn}

\begin{defn}[覆盖结构]
\ollabel{defn:covered-structure}
如果一个结构的论域中的每个元素都是某个闭项的值，则称该结构为\emph{覆盖结构}。
\end{defn}

\begin{ex}
设语言~$\Lang L$ 包含常项符号 $\Obj{zero}$、$\Obj{one}$、$\Obj{two}$、……，
二元谓词符号~$<$，以及二元函数符号 $+$ 和 $\times$。
则 $\Lang L$ 的一个结构~$\Struct M$ 具有论域 $\Domain M = \{0, 1, 2, \ldots
\}$
及赋值 $\Assign{\Obj{zero}}{M} = 0$、$\Assign{\Obj{one}}{M} = 1$、$\Assign{\Obj{two}}{M} = 2$ 等。
对于二元关系符号 $<$，集合 $\Assign{<}{M}$ 是所有满足 $c_1$ 小于~$c_2$ 的有序对 $\tuple{c_1, c_2} \in \Domain{M}^2$ 的集合：
例如，$\tuple{1, 3} \in
\Assign{<}{M}$ 在其中，但 $\tuple{2, 2} \notin \Assign{<}{M}$ 不在。
对于二元函数符号 $+$，按通常方式定义 $\Assign{+}{M}$——例如，$\Assign{+}{M}(2,3)$ 映到~$5$；对二元函数符号~$\times$ 也类似定义。
因此，$\Obj{four}$ 的值就是~$4$，而 $\times(\Obj{two},
+(\Obj{three},\Obj{zero}))$（或用中缀记法，$\Obj{two} \times
(\Obj{three} + \Obj{zero})$）的值是 \begin{multline*}
\Value{\times(\Obj{two}, +(\Obj{three},\Obj{zero}))}{M} =\\
\begin{aligned}
& =\Assign{\times}{M}(\Value{\Obj{two}}{M}, \Value{+(\Obj{three}, \Obj{zero})}{M})\\
& = \Assign{\times}{M}(\Value{\Obj{two}}{M}, \Assign{+}{M}(\Value{\Obj{three}}{M},
\Value{\Obj{zero}}{M})) \\
& = \Assign{\times}{M}(\Assign{\Obj{two}}{M}, \Assign{+}{M}(\Assign{\Obj{three}}{M},
\Assign{\Obj{zero}}{M})) \\
& = \Assign{\times}{M}(2, \Assign{+}{M}(3, 0)) \\
& = \Assign{\times}{M}(2, 3) \\
& = 6
\end{aligned}
\end{multline*}。
\end{ex}

\begin{prob}
算术的标准模型 $\Struct N$ 是覆盖结构吗？请解释。
\end{prob}

\end{document}